\section{Tier 1: Thread-Level Cache (Request-Scoped)}
The thread-level cache is a per-request \texttt{IORef HashMap} that memoises lookups for the lifetime of a single request. It is created at request entry, populated on the first lookup of each key, consulted on subsequent lookups, and discarded at request completion.

\subsection{Correctness Guarantee}
Because the cache is destroyed at request termination and is not shared between threads, it cannot serve stale data: any value it returns is one that the same request has already observed. This is the strongest possible correctness guarantee --- equivalent to local variable assignment.

\subsection{Expected and Realised Impact}
Empirically the thread-level cache achieves a 22.3\% hit rate on eligible lookups, eliminating roughly 2--4 redundant network calls per request.

\section{Tier 2: Server-Local Cache (Node-Scoped)}
The server-local cache is a node-scoped, concurrent cache implemented as a \texttt{TVar} containing a bounded LRU map with TTL-based expiry. It is shared across all request-handling threads on a node.

\subsection{TTL Selection Rationale}
TTLs are configured per data class: 30\,s for routing tables, 60\,s for feature flags, 120\,s for merchant configuration. The choice trades freshness against hit rate; production measurements show that TTLs in this range yield hit rates of 38--47\% with maximum staleness windows that are tolerable for the corresponding data classes (validated by the cache-safety criteria).

\subsection{Realised Impact}
The server-local cache achieves a 41.7\% hit rate on eligible cross-request lookups, eliminating roughly 2{,}800 cross-node calls per second at peak across the cluster.

\section{Cache-Safety Boundaries}
\label{sec:safety}
Not all data is safely cacheable. We define four criteria that a data class must jointly satisfy to be eligible:

\begin{enumerate}[leftmargin=*]
    \item \textbf{Immutability within TTL.} The data must not change over the cache TTL window with probability greater than an agreed bound (typically $10^{-3}$ for configuration classes).
    \item \textbf{Request-Scope Independence.} The value depends only on global or slow-changing state, not on per-request parameters.
    \item \textbf{Semantic Idempotency of Stale Reads.} Even if a stale value is read, the resulting business outcome remains semantically correct (i.e., the worst-case stale read does not produce financial loss).
    \item \textbf{Deterministic Cache Key.} The key is derivable deterministically from the request without collision risk.
\end{enumerate}

A data class is cache-eligible if and only if it satisfies all four criteria. Crucially, \emph{single-resource state} (inventory counts, unique-token availability, account balances on the write path) fails criterion (3): a stale read can lead to double-allocation. Such state is therefore \textbf{ineligible for caching} and is handled by the contention layer described in the next chapter.

% =============================================================================
